\section{Thesis Structure}
\label{sec:intro-structure}

The thesis follows the path from the original circuit to the tested supply. Chapter~\ref{ch:background} explains the conversion principle and the requirements of powering a USB-C device. Chapter~\ref{ch:requirements} turns that application into targets, operating states, and an architecture. Chapter~\ref{ch:analysis} gives the component calculations and explains the interactions between the input and output paths. Chapter~\ref{ch:simulation} then identifies what the supplied models can explain and where a model of the final prototype is still needed.

Chapters~\ref{ch:magnetics} and \ref{ch:hardware} show how the supplied magnetic components were integrated and how the hardware developed from the first prototype through V2 and the reduced V3 layout. Chapter~\ref{ch:control} follows the MCU firmware development, its adaptation for V2, and the operating states used in testing. Chapter~\ref{ch:methods} defines the test setup and measurement methods.

Chapter~\ref{ch:results} brings the implementation and measurements together. Chapter~\ref{ch:discussion} examines the performance limits and the next development steps, and Chapter~\ref{ch:conclusion} answers the research questions. The appendices provide the equations, complete schematics, V2 PCB prints, firmware documentation, and measurement tables needed to follow the work.